\begin{lemma}
\label{lemma-dual-perfect-complex}
Let $(X, \mathcal{O}_X)$ be a ringed space. Let $K$ be a perfect object of
$D(\mathcal{O}_X)$. Then $K^\vee = R\SheafHom(K, \mathcal{O}_X)$ is a
perfect object too and $(K^\vee)^\vee \cong K$. There are
functorial isomorphisms
$$
M \otimes^\mathbf{L}_{\mathcal{O}_X} K^\vee = R\SheafHom(K, M)
$$
and
$$
H^0(X, M \otimes^\mathbf{L}_{\mathcal{O}_X} K^\vee) =
\Hom_{D(\mathcal{O}_X)}(K, M)
$$
for $M$ in $D(\mathcal{O}_X)$.
\end{lemma}

\begin{proof}
By Lemma \ref{lemma-internal-hom-evaluate} there is a canonical map
$$
K = R\SheafHom(\mathcal{O}_X, \mathcal{O}_X)
\otimes_{\mathcal{O}_X}^\mathbf{L} K \longrightarrow
R\SheafHom(R\SheafHom(K, \mathcal{O}_X), \mathcal{O}_X) =
(K^\vee)^\vee
$$
which is an isomorphism by Lemma \ref{lemma-internal-hom-evaluate-isom}.
To check the other statements we will use without further mention that
formation of internal hom commutes with restriction to opens
(Lemma \ref{lemma-restriction-RHom-to-U}).
We may check $K^\vee$ is perfect locally on $X$.
By Lemma \ref{lemma-dual}
to see the final statement it suffices to check that the map
(\ref{equation-eval})
$$
M \otimes^\mathbf{L}_{\mathcal{O}_X} K^\vee
\longrightarrow
R\SheafHom(K, M)
$$
is an isomorphism. This is local on $X$ as well.
Hence it suffices to prove these two statements $K$ is represented
by a strictly perfect complex.

\medskip\noindent
Assume $K$ is represented by the strictly perfect complex
$\mathcal{E}^\bullet$. Then it follows from
Lemma \ref{lemma-Rhom-strictly-perfect}
that $K^\vee$ is represented by the complex whose terms are
$(\mathcal{E}^{-n})^\vee =
\SheafHom_{\mathcal{O}_X}(\mathcal{E}^{-n}, \mathcal{O}_X)$
in degree $n$. Since $\mathcal{E}^{-n}$ is a direct summand of a finite
free $\mathcal{O}_X$-module, so is $(\mathcal{E}^{-n})^\vee$.
Hence $K^\vee$ is represented by a strictly perfect complex too
and we see that $K^\vee$ is perfect.
To see that (\ref{equation-eval}) is an isomorphism, represent
$M$ by a complex $\mathcal{F}^\bullet$.
By Lemma \ref{lemma-Rhom-strictly-perfect} the complex
$R\SheafHom(K, M)$ is represented by the complex with terms
$$
\bigoplus\nolimits_{n = p + q}
\SheafHom_{\mathcal{O}_X}(\mathcal{E}^{-q}, \mathcal{F}^p)
$$
On the other hand, the object $M \otimes^\mathbf{L}_{\mathcal{O}_X} K^\vee$
is represented by the complex with terms
$$
\bigoplus\nolimits_{n = p + q}
\mathcal{F}^p \otimes_{\mathcal{O}_X} (\mathcal{E}^{-q})^\vee
$$
Thus the assertion that (\ref{equation-eval}) is an isomorphism
reduces to the assertion that the canonical map
$$
\mathcal{F}
\otimes_{\mathcal{O}_X}
\SheafHom_{\mathcal{O}_X}(\mathcal{E}, \mathcal{O}_X)
\longrightarrow
\SheafHom_{\mathcal{O}_X}(\mathcal{E}, \mathcal{F})
$$
is an isomorphism when $\mathcal{E}$ is a direct summand of a finite
free $\mathcal{O}_X$-module and $\mathcal{F}$ is any $\mathcal{O}_X$-module.
This follows immediately from the corresponding statement when
$\mathcal{E}$ is finite free.
\end{proof}
